\chapter{Literature Survey}
\label{ch:lit}

This chapter reviews the work on which MemGate builds. It first gives the
background of the overall topic: why LLM agents need memory, how long
contexts fail, and the older ideas from computer systems and psychology that
memory layers borrow. It then surveys the most relevant papers one by one,
compares them in a table, and ends by identifying the research gap that
defines the problem addressed in this project.

\section{Introduction to the overall topic}

\subsection{LLM agents and their memory}
Large language models learn to perform new tasks from instructions and a few
examples placed in the prompt \cite{brown2020gpt3}. Building on this, agent
frameworks interleave reasoning with actions such as tool calls
\cite{yao2023react} and let a model learn from its own failures by writing
verbal reflections into its context \cite{shinn2023reflexion}. Two broad
surveys describe how such agents are assembled from a profile, a planner, a
memory and an action module
\cite{wang2024agentsurvey,xi2023agentsurvey}, and the CoALA framework treats
memory explicitly as working, episodic, semantic and procedural stores around
the model \cite{sumers2024coala}. A survey devoted to agent memory organises
the field by memory source, form and operation, and concludes that deciding
what to remember and what to forget remains the open problem
\cite{zhang2024memorysurvey}.

\subsection{Why longer context windows are not enough}
One response is to make the window larger. YaRN extends the context of
existing models by interpolating rotary position embeddings
\cite{peng2024yarn}. Larger windows, however, are used unevenly. Models make
best use of information at the beginning and the end of a long input and
lose it in the middle \cite{liu2024lostmiddle}; the RULER benchmark finds
effective context lengths far below the advertised ones \cite{hsieh2024ruler};
and a 2025 study of 18 frontier models reports accuracy losses of 30--50\% well
before the documented limit \cite{chroma2025contextrot}. Long-context
benchmarks such as LongBench measure this capacity directly, without a memory
layer \cite{bai2024longbench}.

\subsection{Retrieval as the default baseline}
Retrieval-augmented generation (RAG) keeps external text and puts the passages
most similar to the query into the prompt \cite{lewis2020rag}. Dense passage
retrieval made this practical with learned sentence embeddings
\cite{karpukhin2020dpr}, and RETRO showed that retrieval from a very large
store can stand in for model size \cite{borgeaud2022retro}. Sentence-BERT
\cite{reimers2019sbert} and the distilled MiniLM model \cite{wang2020minilm}
provide the small, fast embeddings that this project uses. RAG stores
everything and never forgets, so it is the natural control against which a
forgetting policy must be judged.

\subsection{Ideas borrowed from systems and psychology}
Memory layers for agents reuse older ideas. The multi-store model of human
memory separates a short-term store from a long-term store with control
processes between them \cite{atkinson1968memory}. Operating systems face the
same choice when physical memory is full: Belady's study of page replacement
\cite{belady1966replacement} and Denning's working-set model
\cite{denning1968workingset} are the classical treatments, and MemGPT makes
the analogy explicit \cite{packer2023memgpt}. Rate--distortion theory
\cite{shannon1959rd,cover2006elements} gives the language for trading the size
of a stored representation against its fidelity, and has recently been
proposed as a framework for agent memory \cite{colaco2026rd}.

\subsection{How memory is evaluated}
LoCoMo provides very long multi-session dialogues with question--answer pairs
whose evidence turns are marked \cite{maharana2024locomo}. LongMemEval
separates several abilities of chat assistants over long histories and
measures the gap between oracle and self-selected evidence
\cite{wu2024longmemeval}. REALTALK applies LoCoMo's annotation scheme to 21 days
of real messages between people \cite{lee2025realtalk}; BEAM scales
conversations to ten million tokens \cite{tavakoli2026beam}; and MemBench
tests factual and reflective memory of agents \cite{tan2025membench}.
LLM-as-a-judge evaluation \cite{zheng2023judge} is widely used for open-ended
answers, while extractive question answering is scored by string match
\cite{rajpurkar2016squad}. For comparing two systems on the same questions,
McNemar's test \cite{mcnemar1947,dietterich1998tests} and the bootstrap
\cite{efron1993bootstrap} are standard; when questions are grouped into
conversations, the resampling must respect that grouping
\cite{cameron2008cluster}. Finally, studies of data leakage show how easily an
evaluation can be contaminated by information that would not be available at
prediction time \cite{kaufman2012leakage,kapoor2023leakage}, a failure this
project met and fixed (Section~\ref{sec:leak}).

\section{Survey of the most relevant work}

This section discusses the most relevant papers in more detail, grouped by
the decision they make.

\subsection{Memory systems for agents}

\textbf{MemGPT} (Packer et al., 2023) \cite{packer2023memgpt} treats the
context window like the main memory of an operating system. The model itself
decides, through function calls, what to page out to an archive and what to
bring back, and a recursive summary keeps a compressed record of evicted
messages. MemGPT shows that an agent can work well beyond its window, and it
was chosen as the base paper of this project. Its limitation, for our
purpose, is that every memory decision is left to the LLM, so there is no
explicit importance policy whose value can be measured on its own.

\textbf{MemoryBank} (Zhong et al., 2024) \cite{zhong2024memorybank} adds a
long-term memory with an Ebbinghaus-style forgetting curve: memories decay
with elapsed time and are reinforced when recalled. It produces human-like
forgetting, but it forgets by age rather than by importance and does not hold
the store to a fixed budget.

\textbf{Generative Agents} (Park et al., 2023) \cite{park2023generative}
record a stream of observations for each simulated character and retrieve
from it by a weighted sum of recency, importance (rated by the LLM) and
relevance, with periodic reflections that summarise the stream. The design
sustains believable long-horizon behaviour, but the stream is never bounded
and the system is not built to reduce tokens or cost.

\textbf{A-MEM} (Xu et al., 2025) \cite{xu2025amem} organises memories as
linked notes in the style of a Zettelkasten. Each new memory is written as a
structured note, linked to related notes, and may cause older notes to be
revised. On LoCoMo it improves multi-hop reasoning while using fewer tokens.
Its focus is on how notes are linked rather than on what to forget under a
strict budget.

\textbf{Mem0} (Chhikara et al., 2025) \cite{chhikara2025mem0} is a
production-oriented memory layer. An LLM extracts candidate facts from each
exchange and chooses among ADD, UPDATE, DELETE and NOOP against the nearest
stored facts. It reports large savings in latency and tokens against
full-context baselines on LoCoMo. Because the extraction and update decisions
are made by an LLM end to end, the contribution of each decision cannot be
separated, which is why Section~\ref{sec:external} rebuilds its memory
structure inside the MemGate harness.

\subsection{Compressing the prompt or the cache}

\textbf{LLMLingua} (Jiang et al., 2023) \cite{jiang2023llmlingua} compresses a
given prompt by removing tokens that a small language model finds predictable,
under a budget controller, reaching up to about 20$\times$ compression with
small losses. It compresses what it is given but does not decide what to
store or forget across sessions.

\textbf{Gist tokens} (Mu et al., 2023) \cite{mu2023gist} train a model to
compress an instruction into a few virtual tokens that can be cached and
reused. \textbf{AutoCompressors} (Chevalier et al., 2023)
\cite{chevalier2023autocompressor} recursively compress long documents into
summary vectors used as soft prompts, and the \textbf{in-context autoencoder}
(Ge et al., 2024) \cite{ge2024icae} packs a context into a small number of
memory slots. All three require training the model and produce opaque
vectors rather than readable memories.

\textbf{StreamingLLM} (Xiao et al., 2024) \cite{xiao2024streaming} observes
that models attend heavily to the first few tokens (``attention sinks'') and
keeps those plus a rolling window of recent tokens in the KV cache, giving
stable generation over millions of tokens. \textbf{H\textsubscript{2}O}
(Zhang et al., 2023) \cite{zhang2023h2o} keeps the ``heavy-hitter'' tokens that
receive the most attention and evicts the rest. Both operate inside one
forward pass over a contiguous context; anything they evict is gone for good.
Compressive Transformers \cite{rae2020compressive}, Memorizing Transformers
\cite{wu2022memorizing} and LongMem \cite{wang2023longmem} go further and
change the model so that it can compress or index its own past activations.

\subsection{Benchmarks and theory}

\textbf{LoCoMo} (Maharana et al., 2024) \cite{maharana2024locomo} contains ten
very long conversations, generated by LLM agents with personas and event
graphs and then checked by people, with about 1{,}500 questions in five
categories whose evidence turns are marked. The marked evidence is what lets
a retrieval-level metric be computed without running a generative model, and
LoCoMo is the main corpus of this project.

\textbf{LongMemEval} (Wu et al., 2024) \cite{wu2024longmemeval} tests five
memory abilities (information extraction, multi-session reasoning, temporal
reasoning, knowledge updates and abstention) over histories of up to about
1.5 million tokens. Its key observation for this project is the 34-point gap
between oracle and self-selected evidence.

\textbf{REALTALK} (Lee et al., 2025) \cite{lee2025realtalk} collects 21 days of
messaging between real people and annotates it in LoCoMo's scheme. Because
LoCoMo is generated by an LLM, REALTALK is the natural test of whether a
scorer learned on LoCoMo has learned what people ask about or only what
generated dialogue looks like.

\textbf{The agent memory survey} (Zhang et al., 2024)
\cite{zhang2024memorysurvey} organises the field and states the open
problem. \textbf{The rate--distortion framework} (Colaco and Lahjouji, 2026)
\cite{colaco2026rd} argues that keeping, compressing and forgetting are one
decision about how much to retain and at what fidelity under a budget, and
predicts that compression should win once the budget is tight enough. It is a
framework rather than a measured system, and Section~\ref{sec:scaling} tests
its prediction.

\section{Comparison of the surveyed work}

Table~\ref{tab:litcompare} compares the main systems along the lines used in
the Review~1 presentation, with the gap each one leaves and how MemGate
addresses it.

{\renewcommand{\arraystretch}{1.15}
\fontsize{9}{11}\selectfont
\begin{longtable}{@{}>{\raggedright\arraybackslash}p{2.3cm}>{\raggedright\arraybackslash}p{2.3cm}>{\raggedright\arraybackslash}p{2.7cm}>{\raggedright\arraybackslash}p{2.6cm}>{\raggedright\arraybackslash}p{2.7cm}@{}}
\caption{Comparison of the main related systems and the gap each leaves.}
\label{tab:litcompare}\\
\toprule
\textbf{Work} & \textbf{Data / benchmark} & \textbf{Technique} & \textbf{Limitation} & \textbf{Addressed in MemGate} \\
\midrule
\endfirsthead
\toprule
\textbf{Work} & \textbf{Data / benchmark} & \textbf{Technique} & \textbf{Limitation} & \textbf{Addressed in MemGate} \\
\midrule
\endhead
MemGPT \cite{packer2023memgpt} (2023) & Multi-Session Chat; document QA & Tiered memory; LLM pages context via function calls & Memory decisions left to the LLM; no measurable importance policy & An explicit, swappable scoring policy, measured on its own \\
MemoryBank \cite{zhong2024memorybank} (2024) & Long-term dialogue (SiliconFriend) & Vector memory with Ebbinghaus forgetting curve & Forgets by elapsed time, not importance; no budget & Forgets by scored utility under a hard storage budget \\
Generative Agents \cite{park2023generative} (2023) & Sandbox simulation & Memory stream; recency, importance and relevance; reflection & Unbounded stream; not built to cut tokens & Bounds the store; keeps retrieval by relevance \\
A-MEM \cite{xu2025amem} (2025) & LoCoMo & Linked notes, dynamic revision & Focus on linking, not on budgeted forgetting & Budget-driven forgetting, compared on LoCoMo \\
Mem0 \cite{chhikara2025mem0} (2025) & LoCoMo & LLM fact extraction with ADD, UPDATE and DELETE & Decisions made end to end by an LLM & Rebuilt in the same harness and compared at equal budgets \\
StreamingLLM \cite{xiao2024streaming} (2024) & Long-stream language modelling & Attention sinks plus rolling KV cache & Evicted content is gone for good & Evicted items can live on as gists or facts and be recalled \\
H\textsubscript{2}O \cite{zhang2023h2o} (2023) & OPT and LLaMA tasks & Keeps heavy-hitter tokens in the KV cache & Works inside one forward pass & Works across sessions, at the conversation level \\
LLMLingua \cite{jiang2023llmlingua} (2023) & GSM8K, BBH and others & Token-level prompt compression & Compresses, but does not decide what to keep & Tests compression against dropping inside one policy \\
Gist tokens \cite{mu2023gist} (2023) & Instruction following & Learned virtual tokens & Needs training; fixed instructions only & No training of the base model; readable memory \\
AutoCompressors \cite{chevalier2023autocompressor} (2023) & Long-context language modelling & Recursive summary vectors & Opaque vectors; no explicit forgetting & Readable gists plus a retrievable fact store \\
Memory survey \cite{zhang2024memorysurvey} (2024) & Survey & Taxonomy of agent memory & No system & Measures several of its open questions \\
Rate--distortion \cite{colaco2026rd} (2026) & Theory & Memory as rate--distortion & Framework, not a measured system & Tests its crossover prediction to 182$\times$ \\
\bottomrule
\end{longtable}}

\section{Research gap and problem definition}
\label{sec:gap}

The survey leads to four observations.
\begin{enumerate}
\item \textbf{Tiered memory with forgetting is no longer new.} MemGPT, Mem0,
  A-MEM and others already ship it. A project that only builds another
  three-tier layer would repeat their work.
\item \textbf{No one has measured the decisions separately.} Each system is
  evaluated as a whole, with its own model, prompts and budgets, so a reader
  cannot tell whether its gains come from what it keeps, how it compresses,
  how it retrieves, or the model it uses.
\item \textbf{The accounting is incomplete.} Evaluations cap the context sent
  to the model but leave the store unbounded, which lets a store-everything
  baseline keep its whole history for free. Byte costs, including the
  embedding index, are not counted.
\item \textbf{Salience scoring is ad hoc.} Importance is rated by hand-written
  rules or by asking an LLM, and whether a learned scorer would do better, and
  whether it would carry over from generated to human conversation, has not
  been tested.
\end{enumerate}

These observations define the problem stated in Section~1.3: build a memory
layer whose keep, compress, forget and recall decisions can be isolated and
swapped; hold it to a storage budget counted in tokens and in bytes; and
measure each decision, including a learned scorer, against controls that
differ from it in one factor only, with statistics that respect the grouping
of questions into conversations.

\section*{Summary}
The literature offers many working memory systems, several ways to compress
prompts and caches, and good long-conversation benchmarks. What it lacks is a
controlled measurement of the individual memory decisions under a bounded
store, a byte-level accounting of their cost, and a tested learned scorer.
MemGate is designed to fill exactly that gap.
